\documentclass{article}
\usepackage{geometry}
\geometry{margin=1.5cm, vmargin={0pt,1cm}}
\setlength{\topmargin}{-1cm}
\setlength{\headheight}{12.5pt}
\setlength{\paperheight}{29.7cm}
\setlength{\textheight}{25.3cm}

% useful packages.
\usepackage[UTF8]{ctex}
\usepackage{fancyhdr}
\usepackage{layout}
\usepackage{luacode}

% import common macros
% % % % % % % % % % % % % % % % % % % % % % % % % % % % % % % %
% Common Macros for LaTeX
% % % % % % % % % % % % % % % % % % % % % % % % % % % % % % % %

% Useful packages
\usepackage{amsfonts}
\usepackage{amsmath}
\usepackage{amssymb}
\usepackage{amsthm}
\usepackage{enumerate}
\usepackage{graphicx}
\usepackage{multicol}
\usepackage[ruled,linesnumbered]{algorithm2e}

% Math operators and common symbols
\newcommand{\dif}{\mathrm{d}}
\newcommand{\innerProd}[1]{\left\langle #1 \right\rangle}
\newcommand{\difFrac}[2]{\frac{\dif #1}{\dif #2}}
\newcommand{\pdfFrac}[2]{\frac{\partial #1}{\partial #2}}
\newcommand{\T}{\mathrm{T}}
\newcommand{\norm}[1]{\left\| #1 \right\|}
\newcommand{\abs}[1]{\left| #1 \right|}

% Vector and Matrix shortcuts
\newcommand{\vecTwo}[2]{
  \begin{bmatrix}
    #1 \\
    #2
\end{bmatrix}}
\newcommand{\vecThree}[3]{
  \begin{bmatrix}
    #1 \\
    #2 \\
    #3
\end{bmatrix}}
\newcommand{\vecTwoH}[2]{
  \begin{bmatrix}
    #1 & #2
  \end{bmatrix}
}
\newcommand{\vecThreeH}[3]{
  \begin{bmatrix}
    #1 & #2 & #3
  \end{bmatrix}
}

% Common fields
\newcommand{\R}{\mathbb{R}}
\newcommand{\C}{\mathbb{C}}
\newcommand{\Z}{\mathbb{Z}}
\newcommand{\N}{\mathbb{N}}

% Solutions and proofs
\newenvironment{solution}{
\begin{proof}[解]}{
\end{proof}}

% Utility to get environment variables (requires lualatex)
\newcommand{\getEnv}[2]{\directlua{tex.print(os.getenv("#1") or "#2")}}


\usepackage{listings}
\usepackage{xcolor}

% Configure listings for Python code highlighting
\lstset{
  language=Python,
  basicstyle=\ttfamily\small,
  keywordstyle=\bfseries\color{blue},
  commentstyle=\color{gray},
  stringstyle=\color{red},
  breaklines=true,
  showstringspaces=false,
  tabsize=4,
  frame=single,
  framesep=5pt,
  rulecolor=\color{gray},
  backgroundcolor=\color{white},
  captionpos=b,
  numberstyle=\tiny\color{gray},
  numbers=left,
  stepnumber=5,
}

\lstnewenvironment{plaintext}[1][]{
  \lstset{
    language=none,
    basicstyle=\ttfamily,
    keywordstyle=,
    commentstyle=,
    stringstyle=,
    numbers=none,
    frame=none,
    backgroundcolor=\color{white},
    #1                % 允许用户自行覆盖默认设置
  }
}{}

% Name and uID
\newcommand{\name}{\getEnv{NAME}{NAME}}
\newcommand{\uid}{\getEnv{UID}{UID}}
\newcommand{\course}{数值代数}
\newcommand{\hwNum}{11}

\newcommand{\fl}{\mathrm{fl}}

\begin{document}

\pagestyle{fancy}
\fancyhead{}
\lhead{\name (\uid)}
\chead{\course \#\hwNum}
\rhead{\today}

\section{题目 1}

设 $x_k$ 是由最速下降法产生的。证明
\[
  \phi(x_{k+1}) + b^T A^{-1} b \leq \left(1 - \frac{1}{\kappa_2(A)}\right)
  \left(\phi(x_k) + b^T A^{-1} b\right),
\]
其中 $\kappa_2(A) = \norm{A}_2 \norm{A^{-1}}_2$。

\begin{solution}
  考虑：
  \[
    \nabla \phi(x) = 2 Ax - 2b = 0
  \]
  令 $x^* = A^{-1} b$，则 $\nabla \phi(x^*) = 0$。代入得：
  \[
    \phi(x^*) = b^T A^{-1} b - 2b^T A^{-1}b = - b^T A^{-1} b
  \]
  于是有：
  \[
    \phi(x_k) + b^T A^{-1} b = \phi(x_k) - \phi(x^*)
  \]
  令 $e_k = x_k - x^*$，则：
  \[
    \phi(x_k) - \phi(x^*) = (e_k)^T A e_k.
  \]
  于是只需证明：
  \[
    \phi(x_k) - \phi(x_{k+1}) \geq \frac{1}{\kappa_2(A)}
    (\phi(x_k) - \phi(x^*)) = \frac{1}{\|A\| \|A^{-1}\|}
    (e_k)^T A e_k.
  \]
  令 $r_k = b - Ax_k = -Ae_k$ 于是有：
  \[
    \phi(x_k) - \phi(x_{k+1}) =
    \frac{(r_k^Tr_k)^2}{r_k^TAr_k}=
    \frac{\|r_k\|^4}{r_k^TAr_k}
  \]
  于是只需证明：
  \[
    \|r_k\|^4 \|A\|\|A^{-1}\| \geq (r_k^T A r_k) (e_k^TA e_k)
  \]
  又由于 $e_k = -A^{-1}r_k$，于是有：
  \[
    e_k^T A e_k = r_k^T A^{-1} r_k
  \]
  另一方面，$r^T A r \leq \|r\| \|Ar\| \leq \|r\|^2 \|A\|$ 由
  Cauchy-Schwarz 不等式和矩阵范数性质易证，从而原不等式成立。

\end{solution}

\section{题目 2}

在共轭梯度法中，证明
\[
  (r_k, r_k) = -\alpha_{k-1} (r_k, A p_{k-1}),
\]
\[
  (r_k, r_k) = \alpha_k (p_k, A p_k).
\]

\begin{solution}

  \[
    \begin{aligned}
      r^T_k r_k &= r^T_k (r_{k-1} - \alpha_{k-1}Ap_{k-1}) = r^T_k
      r_{k-1} - \alpha_{k-1} r^T_k A p_{k-1} = -\alpha_{k-1} r^T_k A p_{k-1} \\
      &= \frac{r^T_{k} r_{k}}{p^T_{k} A p_{k}} p^T_{k} A p_{k} =
      \alpha_k p^T_k A p_k.
    \end{aligned}
  \]

\end{solution}

\section{题目 3}

设
\[
  A =
  \begin{bmatrix} 2 & 1 & 0 \\ 1 & 2 & 1 \\ 0 & 1 & 1
  \end{bmatrix}, \quad
  b =
  \begin{bmatrix} 1 \\ 0 \\ 1
  \end{bmatrix}.
\]
请用共轭梯度法求 $Ax = b$。

\begin{solution}
  由顺序主子式 $2 > 0$，$\det\matTwoTwo{2}{1}{1}{2} = 3 > 0$，$\det A = 1 > 0$
  知 $A$ 对称正定。取 $x_0 = (0,0,0)^T$，则 $r_0 = b - Ax_0 = (1,0,1)^T$，
  $p_0 = r_0 = (1,0,1)^T$。

  \textbf{第 1 步：}
  \[
    Ap_0 = \vecThree{2}{2}{1}, \quad r_0^T r_0 = 2, \quad p_0^T A p_0 = 3,
    \quad \alpha_0 = \frac{2}{3}.
  \]
  \[
    x_1 = x_0 + \alpha_0 p_0 = \vecThree{2/3}{0}{2/3}, \quad
    r_1 = r_0 - \alpha_0 Ap_0 = \vecThree{-1/3}{-4/3}{1/3}.
  \]
  \[
    r_1^T r_1 = \frac{18}{9} = 2, \quad
    \beta_0 = \frac{r_1^T r_1}{r_0^T r_0} = 1, \quad
    p_1 = r_1 + \beta_0 p_0 = \vecThree{2/3}{-4/3}{4/3}.
  \]

  \textbf{第 2 步：}
  \[
    Ap_1 = \vecThree{0}{-2/3}{0}, \quad p_1^T Ap_1 = \frac{8}{9}, \quad
    \alpha_1 = \frac{2}{8/9} = \frac{9}{4}.
  \]
  \[
    x_2 = x_1 + \alpha_1 p_1 = \vecThree{2/3+3/2}{-3}{2/3+3}
    = \vecThree{13/6}{-3}{11/3}.
  \]
  \[
    r_2 = r_1 - \alpha_1 Ap_1 = \vecThree{-1/3}{-4/3+3/2}{1/3}
    = \vecThree{-1/3}{1/6}{1/3}.
  \]
  \[
    r_2^T r_2 = \frac{1}{4}, \quad
    \beta_1 = \frac{1/4}{2} = \frac{1}{8}, \quad
    p_2 = r_2 + \beta_1 p_1 = \vecThree{-1/4}{0}{1/2}.
  \]

  \textbf{第 3 步：}
  \[
    Ap_2 = \vecThree{-1/2}{1/4}{1/2}, \quad p_2^T Ap_2 =
    \frac{1}{8}+\frac{1}{4} =
    \frac{3}{8}, \quad \alpha_2 = \frac{1/4}{3/8} = \frac{2}{3}.
  \]
  \[
    x_3 = x_2 + \alpha_2 p_2 = \vecThree{13/6-1/6}{-3}{11/3+1/3}
    = \vecThree{2}{-3}{4}.
  \]

  故 $Ax = b$ 的解为 $x = (2, -3, 4)^T$。
\end{solution}

\section{题目 4}

$A$ 为 $n$ 阶对称正定矩阵，$p_1, p_2, \dots, p_n$ 为其共轭向量系。证明
\[
  A^{-1} = \sum_{k=1}^{n} \frac{p_k p_k^T}{p_k^T A p_k}.
\]

\begin{solution}
  由于 $p_1, \dots, p_n$ 为 $A$-共轭非零向量组，它们线性无关，构成 $\R^n$ 的一组基。
  对任意 $v \in \R^n$，设 $w = A^{-1}v$，即 $Aw = v$。将 $w$ 在此基下展开：
  \[
    w = \sum_{k=1}^n c_k p_k.
  \]
  两边左乘 $A$ 得 $v = Aw = \sum_k c_k Ap_k$。由 $A$-共轭性，
  \[
    p_k^T v = p_k^T \sum_j c_j Ap_j = c_k \, p_k^T Ap_k,
  \]
  故 $c_k = \dfrac{p_k^T v}{p_k^T A p_k}$。于是
  \[
    A^{-1}v = w = \sum_{k=1}^n \frac{p_k^T v}{p_k^T A p_k} \, p_k
    = \sum_{k=1}^n \frac{p_k p_k^T}{p_k^T A p_k} \, v.
  \]
  上式对任意 $v \in \R^n$ 成立，因此
  \[
    A^{-1} = \sum_{k=1}^n \frac{p_k p_k^T}{p_k^T A p_k}.
  \]
\end{solution}

\section{题目 5}

试证：当最速下降法在有限步求得极小值时，最后一步迭代的下降方向必是 $A$ 的一个特征向量。

\begin{solution}

  \[
    x_{k+1} = x_k + \alpha_{k}p_{k} = x^* = A^{-1} b
  \]
  考虑梯度 $\nabla \phi(x_{k+1}) = 0$，于是有：
  \[
    Ax_k+\alpha_k A p_k - b = 0
  \]
  又有 $p_k = b-Ax_k$，从而有：
  \[
    \alpha A p_k = b-Ax_k = p_k \Rightarrow A p_k = \frac{1}{\alpha} p_k,
  \]
  即 $p_k$ 是 $A$ 的特征向量。

\end{solution}

\section{题目 6}

设 $A$ 是一个只有 $k$ 个互不相同特征值的 $n \times n$ 实对称矩阵，$r$ 是任一 $n$ 维实向量。

证明：子空间 $\mathrm{span}\{r, Ar, \dots, A^{n-1}r\}$ 的维数至多是 $k$。

\begin{solution}
  由于 $A$ 是实对称矩阵，由谱定理知 $A$ 可对角化：$A = Q \Lambda Q^T$，其中
  $\Lambda = \mathrm{diag}(\lambda_1, \dots, \lambda_n)$，$\lambda_i \in
  \{\mu_1, \dots, \mu_k\}$。$A$ 的最小多项式为
  \[
    m_A(t) = \prod_{j=1}^k (t - \mu_j),
  \]
  这是因为实对称矩阵可对角化，其最小多项式无重根，且包含所有互不相同特征值对应的因子。于是
  \[
    m_A(A) = \prod_{j=1}^k (A - \mu_j I) = 0,
  \]
  即 $A^k$ 可表示为 $I, A, \dots, A^{k-1}$ 的线性组合。对任意向量 $r$，有
  \[
    A^k r \in \mathrm{span}\{r, Ar, \dots, A^{k-1}r\}.
  \]
  由归纳法，$A^j r \in \mathrm{span}\{r, Ar, \dots, A^{k-1}r\}$ 对所有 $j \geq k$ 成立。因此
  \[
    \mathrm{span}\{r, Ar, \dots, A^{n-1}r\} = \mathrm{span}\{r, Ar,
    \dots, A^{k-1}r\},
  \]
  至多由 $k$ 个向量张成，维数至多为 $k$。
\end{solution}

\section{题目 7}

试证：如果系数矩阵 $A$ 至多有 $l$ 个互不相同的特征值，则共轭梯度法至多 $l$ 步就可得到方程组 $Ax = b$ 的精确解。

\begin{solution}
  设 $A$ 为 $n$ 阶对称正定矩阵，$l$ 个互不相同的特征值为 $\mu_1, \dots, \mu_l$。
  取初始迭代 $x_0$，令 $r_0 = b - Ax_0$，$e_0 = x_0 - x^*$。

  在共轭梯度法中，误差满足 $e_k = p_k(A) e_0$，其中 $p_k$ 是满足 $p_k(0) = 1$ 的 $k$
  次多项式。由谱分解 $A = Q\Lambda Q^T$，令 $\hat{e}_0 = Q^T e_0$，则
  \[
    \norm{e_k}_A^2 = e_k^T A e_k = \sum_{i=1}^n \lambda_i |p_k(\lambda_i)|^2
    |\hat{e}_{0,i}|^2.
  \]
  构造 $l$ 次多项式
  \[
    p_l(t) = \prod_{j=1}^l \left(1 - \frac{t}{\mu_j}\right),
  \]
  则 $p_l(0) = 1$，且 $p_l(\mu_j) = 0$（$j = 1, \dots, l$）。由于每个特征值
  $\lambda_i$ 等于某个 $\mu_j$，故 $p_l(\lambda_i) = 0$，从而
  \[
    \norm{e_l}_A^2 = \sum_{i=1}^n \lambda_i \cdot 0 \cdot |\hat{e}_{0,i}|^2 = 0.
  \]
  因此 $e_l = 0$，即 $x_l = x^*$，共轭梯度法至多 $l$ 步得到精确解。
\end{solution}

\section{上机习题}

（教材第 160 页）

（2）用 Hilbert 矩阵测试你所编写的共轭梯度法程序：
\[
  a_{ij} = \frac{1}{i+j+1}, \quad b_i = \frac{1}{3} \sum_{j=1}^{n}
  a_{ij}, \quad 1 \leq i, j \leq n.
\]

（3）分别用 Jacobi 迭代法、G-S 迭代法和共轭梯度法求解下述方程：
\[
  \begin{bmatrix}
    0.3 & 0.1 & 0 & 0 & 0.2 \\
    0.1 & 0.3 & 0.1 & 0 & 0 \\
    0 & 0.1 & 0.3 & 0.1 & 0 \\
    0 & 0 & 0.1 & 0.3 & 0.1 \\
    0.2 & 0 & 0 & 0.1 & 0.3
  \end{bmatrix}
  \begin{bmatrix} x_1 \\ x_2 \\ x_3 \\ x_4 \\ x_5
  \end{bmatrix}
  =
  \begin{bmatrix} 0.6 \\ 0.5 \\ 0.5 \\ 0.5 \\ 0.6
  \end{bmatrix}.
\]

\begin{solution}
  第 (2) 题中，由 $b_i = \frac{1}{3}\sum_j a_{ij}$ 可知精确解为 $x_j =
  \frac{1}{3}$（$j = 1, \dots, n$）。
  共轭梯度法实现见 \lstinline{conjugate_gradient}，停机准则为
  $\norm{x_{k+1} - x_k}_\infty < 10^{-10}$。

  \lstinputlisting[language=Python, firstline=5, lastline=20,
  caption=共轭梯度法]{exercise.py}
  \lstinputlisting[language=Python, firstline=52, lastline=63,
  caption=Hilbert 矩阵测试]{exercise.py}
  \lstinputlisting[language={},caption=第 (2) 题输出]{output.txt}

  $n$ 较小时 CG 在 $n$ 步以内即可收敛到高精度解。随着 $n$ 增大，Hilbert 矩阵的条件数急剧增长（$n=10$
  时已超过 $10^{14}$），
  此时 CG 的收敛精度受浮点误差限制，但迭代次数仍远小于 $n$。

  \bigskip

  第 (3) 题中，该矩阵行和为 $0.5 < 1$，严格对角占优，因此 Jacobi 和 G-S 迭代均收敛。
  精确解为 $x = (1, 1, 1, 1, 1)^T$。

  \lstinputlisting[language=Python, firstline=68, lastline=91,
  caption=5$\times$5 方程组求解]{exercise.py}

  从输出可以看到，CG 仅需 4 步即达到机器精度，而 Jacobi 需要 153 步，G-S 需要 51 步。
  CG 对 $5\times 5$ 的正定系统最多需要 5 步，此处仅用 4 步即达到精确解。
\end{solution}

\end{document}
